\documentclass[10pt]{article}

\usepackage[T1]{fontenc}
\usepackage[utf8]{inputenc}
\usepackage{babel}
\babelprovide[import, main]{polish}
\usepackage[a4paper, margin=1.5cm]{geometry}
\usepackage{graphicx}
\usepackage{listings}
\usepackage{xcolor}
\usepackage{graphicx}

\lstdefinestyle{mystyle}{
    basicstyle=\ttfamily\small,
    frame=single,
    breaklines=true,
    captionpos=b,
    showstringspaces=false
}

\title{Projekt I przetwarzanie równoległe 2025/26 SS, OpenMP, Vtune}
\author{Emilia Bonk 160108 \\ Szymon Urban 160271}
\date{}

\begin{document}

\maketitle

\section*{Wstęp.}

\begin{itemize}
\item \textbf{Nazwa przedmiotu:} Laboratorium z przetwarzania równoległego
\item \textbf{Autorzy:} Emilia Bonk (160108), Szymon Urban (160271)
\item \textbf{Grupa dziekańska:} L09
\item \textbf{Termin zajęć laboratoryjnych:} Wtorek 9:45-11:15
\item \textbf{Termin oddania sprawozdania:} 5.05.2026
\item \textbf{Wersja sprawozdania:} pierwsza
\item \textbf{Opis zadania:}
Celem ćwiczenia była analiza implementacji algorytmu Sita Eratostenesa 
w środowisku równoległym i sekwencyjnym z wykorzystaniem biblioteki OpenMP. 
Zadanie polegało na analizie i modyfikacji podanych kodów, z róznymi 
implementacjami algorytmu, a następnie porównaniu ich wydajności.
\item \textbf{Adres e-mail:} emilia.bonk@student.put.poznan.pl, szymon.urban@student.put.poznan.pl
\end{itemize}

\section{Informacje techniczne.}

\begin{table}[h!]
\centering
\begin{tabular}{|l|l|}
\hline
\textbf{Parametr} & \textbf{Wartość} \\
\hline
Model procesora & 17th Gen Intel(R) Core(TM) i7-12700F \\
\hline
Liczba procesorów fizycznych & 1 \\
\hline
Liczba rdzeni & 12 \\
\hline
Liczba procesorów logicznych & 20 \\
\hline
Pamięć cache L1 & 1,0 MB \\
\hline
Pamięć cache L2 & 12,0 MB \\
\hline
Pamięć cache L3 & 25,0 MB \\
\hline
System operacyjny & Windows 10 \\
\hline
Kompilator & Intel(R) oneAPI DPC++/C++ Compiler 2025.3.3 \\
\hline
Wersja kompilatora & 2025.3.3.20260319 \\
\hline
Program testujący & Intel VTune Profiler 2026.0.0\\
\hline
\end{tabular}
\caption{Informacje techniczne dotyczące środowiska testowego}
\end{table}


\newpage
\section{Prezentacja i analiza kodów.}

\subsection*{K1 - Algorytm sekwencyjny oparty na próbnych dzieleniach.}
\textbf{Oznaczenie skrótowe:} SEQ-DIV-BASIC
\\ \\
Algorytm sekwencyjny oparty na metodzie dzielenia próbnego. Najpierw wyznacza liczby pierwsze w zakresie do $\sqrt{n}$, a następnie szeregowo sprawdza podzielność każdej liczby $i \in [m, n]$ przez te dzielniki. Koszt obliczeń jest zmienny – najniższy dla liczb parzystych, najwyższy dla liczb pierwszych.
\\

\begin{lstlisting}[style=mystyle, caption={Kod programu k1.cpp}, label={lst:k1}]
int calculate(int n, int m){
	bool* result = (bool*)malloc((n - m + 1) * sizeof(bool));
	memset(result, true, (n - m + 1) * sizeof(bool));
	int prime_counter = 0;

	bool* primeArray = (bool*)malloc((sqrt(n) + 1) * sizeof(bool));
	memset(primeArray, true, (sqrt(n) + 1) * sizeof(bool));
	
	if (m <= 0 && n >= 0) result[0 - m] = false;
	if (m <= 1 && n >= 1) result[1 - m] = false;

	for (int i = 2; i * i <= n; i++) {
		for (int j = 2; j * j <= i; j++) {
			if (primeArray[j] == true && i % j == 0) { 
				primeArray[i] = false; 
				break;
			}
		}
	}
	for (int i = m; i <= n; i++){
		for (int j = 2; j * j <= i; j++){
			if (primeArray[j] == true && i % j == 0) {
				result[i - m] = false;
				break; 
			} 
		}
		if(result[i-m] == true){
			prime_counter++;
		}
	}
	free(result);
	free(primeArray);

	return prime_counter;
}
\end{lstlisting}

\newpage

\subsection*{K2 - Liczby pierwsze wyznaczane równolegle przez dzielenie w zakresie <m,n>.}
\textbf{Oznaczenie skrótowe:} PAR-DIV-GUIDED
\\ \\
Wariant równoległy wykorzystujący dzielenie próbne z dyrektywą \texttt{schedule(guided)}. Zastosowanie malejących paczek zadań pozwala na 
efektywne zrównoważenie obciążenia procesora, niwelując różnice w czasie sprawdzania małych i dużych liczb. Brak wyścigów przy zapisie do tablicy wyników.
\\
\begin{lstlisting}[style=mystyle, caption={Kod programu k2.cpp}, label={lst:k2}]
int calculate(int n, int m){
	int prime_counter = 0;

	//...liczenie sita (PrimeArray) jak w kodzie k1
	#pragma omp parallel
	{
	#pragma omp for schedule(guided) reduction(+:prime_counter)
	for (int i = m; i <= n; i++){
		for (int j = 2; j * j <= i; j++){
			if (primeArray[j] == true && i % j == 0) { 
				result[i - m] = false;
				break;
			}
		}	
		if(result[i-m] == true){
			prime_counter++;
		}
	}
	} 
	free(result);
	free(primeArray);
	
	return prime_counter;
}
\end{lstlisting}

\newpage

\subsection*{K3 - Sito sekwencyjne bez lokalności dostępu do danych.}
\textbf{Oznaczenie skrótowe:} SEQ-SIEVE-NL
\\ \\
Sito Eratostenesa realizowane sekwencyjnie. Zamiast dzielenia, algorytm wykreśla wielokrotności liczb pierwszych $i \le \sqrt{n}$ w całym zakresie $[m, n]$. 
Charakteryzuje się niską złożonością obliczeniową $O(n \log \log n)$, ale ograniczoną lokalnością danych przy skokowym dostępie do tablicy wyników.
\\

\begin{lstlisting}[style=mystyle, caption={Kod programu k3.cpp}, label={lst:k3}]
int calculate(int n, int m){
	//...liczenie sita jak we wczesniejszych kodach 

	for (int i = 2; i*i <= n; i++){
		if (primeArray[i]) {
			int firstMultiple = (m / i);
			if (firstMultiple <= 1) {
				firstMultiple = i + i;
			} else if (m % i) { 
				firstMultiple = (firstMultiple * i) + i;
			} else {
				firstMultiple = (firstMultiple * i);
			}
			for (int j = firstMultiple; j <= n; j+=i) { 
				result[j-m] = false;
			}
		}
	}
	if (m <= 0 && n >= 0) result[0 - m] = false;
	if (m <= 1 && n >= 1) result[1 - m] = false;

	for (int i = m; i <= n; i++) {
    	if (result[i - m]) prime_counter++;
	}

	free(result);
	free(primeArray);
	
	return prime_counter;
}
\end{lstlisting}

\newpage
\subsection*{K3a - Sito sekwencyjne z lokalnością dostępu do danych.}
\textbf{Oznaczenie skrótowe:} SEQ-SIEVE-LL
\\ \\ 
Sekwencyjne sito segmentowe zoptymalizowane pod kątem lokalności danych. Zakres $[m, n]$ dzielony jest na bloki mieszczące się w pamięci cache L1/L2.
Przetwarzanie blokowe znacząco przyspiesza operacje zapisu do tablicy wyników poprzez redukcję pudeł cache.
\\

\begin{lstlisting}[style=mystyle, caption={Kod programu k3a.cpp}, label={lst:k3a}]
int calculate(int n, int m){
    //...alokacja pamieci i liczenie PrimeArray tak samo jak w kodzie k3

    int blockSize = 65536; 
    int numberOfBlocks = (n - m) / blockSize;
    if ((n - m) % blockSize != 0) { 
        numberOfBlocks++;
    }

    for (int i = 0; i < numberOfBlocks; i++) {
        int low = m + i * blockSize; 
        int high = low + blockSize - 1;
        if (high > n) {
            high = n;
        }

        for (int j = 2; j * j <= high; j++) {
            if (primeArray[j] == true) {
                int firstMultiple = (low / j);
                if (firstMultiple <= 1) {
                    firstMultiple = j + j;
                } else if (low % j) {
                    firstMultiple = (firstMultiple * j) + j;
                }  else {
                    firstMultiple = (firstMultiple * j);
                }
                for (int k = firstMultiple; k <= high; k += j) {
                    result[k - m] = false;
                }
            }
        }
    }

	//...zliczanie prime_count i zwalnianie pamieci jak w kodzie k3
}
\end{lstlisting}

\newpage
\subsection*{K4 - Sito równoległe funkcyjne bez lokalności dostępu do danych.}
\textbf{Oznaczenie skrótowe:} PAR-SIEVE-NL
\\ \\
Równoległe sito z podziałem funkcyjnym (\texttt{schedule(dynamic, 1)}). Wątki pobierają kolejne liczby pierwsze i wykreślają ich wielokrotności w 
całym zakresie wyników. Występuje zjawisko \texttt{false sharing} oraz bezpieczne wyścigi przy zapisie tej samej wartości logicznej pod ten sam indeks.
\\

\begin{lstlisting}[style=mystyle, caption={Kod programu k4.cpp}, label={lst:k4}]
int calculate(int n, int m){
	//...alokacja pamieci i lizcenie PrimeArray jak wczesniej 

	int limit = sqrt(n);
	#pragma omp parallel for schedule(dynamic, 1)
	for (int i = 2; i <= limit; i++){
		if (primeArray[i]) {
			int firstMultiple = (m / i);
			if (firstMultiple <= 1) {
				firstMultiple = i + i;
			} else if (m % i){ 
				firstMultiple = (firstMultiple * i) + i;
			} else {
				firstMultiple = (firstMultiple * i);
			}
			for (int j = firstMultiple; j <= n; j+=i) {
				result[j-m] = false;
			}
		}
	}

	//zliczanie prime_count i zwalnianie pamieci jak wczesniej
}
\end{lstlisting}

\newpage
\subsection*{K4a - Sito równoległe funkcyjne bez lokalności dostępu do danych z usunięciem redundancji zapisów.}
\textbf{Oznaczenie skrótowe:} PAR-SIEVE-NL-A
\\ \\ 
Modyfikacja algorytmu K4 wprowadzająca warunkowy zapis do tablicy wyników (\texttt{if (result[j-m])}). Ma to na celu ograniczenie liczby operacji zapisu i związanych 
z nimi konfliktów w pamięci cache, co poprawia wydajność przy zachowaniu podziału funkcyjnego.
\\

\begin{lstlisting}[style=mystyle, caption={Kod programu k4a.cpp}, label={lst:k4a}]
int calculate(int n, int m){
	//...alokacja pamieci i zliczanie PrimeArray jak wczesniej

	int limit = sqrt(n);
	#pragma omp parallel for schedule(dynamic, 1)
	for (int i = 2; i <= limit; i++){
		if (primeArray[i]) {
			int firstMultiple = (m / i);
			if (firstMultiple <= 1) {
				firstMultiple = i + i;
			} else if (m % i){ 
				firstMultiple = (firstMultiple * i) + i;
			} else {
				firstMultiple = (firstMultiple * i);
			}
			for (int j = firstMultiple; j <= n; j+=i) {
				if (result[j-m]) result[j-m] = false;
			} 
		}
	}

	//zliczanie prime_counter i zwalnianie pamieci jak w k4
}
\end{lstlisting}


\newpage
\subsection*{K5 - Sito równoległe domenowe z lokalnością dostępu do danych}
\textbf{Oznaczenie skrótowe:} PAR-SIEVE-DOM-LL
\\ \\
Równoległe sito segmentowe z podziałem domenowym. Każdy wątek przetwarza niezależne bloki danych mieszczące się w cache, co eliminuje \texttt{false sharing} i wyścigi. 
Choć wydawałoby się, że \texttt{schedule(static)} minimalizuje narzut synchronizacji przy wyrównanym obciążeniu bloków,
to jednak znacznie lepszy performormance zaobserwowaliśmy przy \texttt{schedule(guided)}.
\\

\begin{lstlisting}[style=mystyle, caption={Kod programu k5.cpp}, label={lst:k5}]
int calculate(m, n){
	//...alokacja pamieci i liczenie PrimeArray jak wczesniej

	int blockSize = 65536; 
	int numberOfBlocks = (n - m) / blockSize + ((n - m) % blockSize != 0);

	#pragma omp parallel for schedule(guided)
	for (int i = 0; i < numberOfBlocks; i++) {
		int low = m + i * blockSize; 
		int high = low + blockSize - 1;
		if (high > n) high = n;

		for (int j = 2; j * j <= high; j++) {
			if (primeArray[j]) {
				int firstMultiple = (low / j) * j;
				if (firstMultiple < low) firstMultiple += j;
				if (firstMultiple == j) firstMultiple += j; 

				for (int k = firstMultiple; k <= high; k += j) {
					result[k - m] = false;
				}
			}
		}
	}

	//zliczanie prime_counter i zwalnianie pamieci jak wczesniej
}
\end{lstlisting}


\newpage
\section{Prezentacja wyników i omówienie przebiegu eksperymentu.}
\subsection*{Zasady pracy z oprogramowaniem VTune.}
\subsection*{Porównanie podstawowych parametrów wywołania dla kodów k1-k5.}
Porównania zostały wykonane dla kodów przedstawionych w poprzedniej sekcji, z tymi samymi schedulingami
i rozmiarami bloków na zakresie liczb od 1 do 100 milionów.
\begin{table}[h!]
\centering
\begingroup
% \resizebox{\textwidth}{!}{%
\begin{tabular}{|l|c|c|c|c|c|c|c|}
\hline
 & \textbf{k1} & \textbf{k2} & \textbf{k3} & \textbf{k3a} & \textbf{k4} & \textbf{k4a} & \textbf{k5} \\
\hline
Elapsed Time [s] & 23.750 & 2.394 & 0.660 & 0.130 & 0.381 & 0.296 & 0.060 \\
\hline
Prędkość przetwarzania [mln l. pierwszych/s] & 4.210 & 2.406 & 14.394 & 35.714 & 15.121 & 19.464 & 96.024 \\
\hline
Instructions retired [mld] & 405,744 & 405,599 & 0,515 & 0,777 & 2,179 & 1,501 & 1,702 \\
\hline
Clockticks [miliardów tików] & 105,112 & 166,746 & 2,914 & 0,566 & 19,808 & 14,777 & 1,995 \\
\hline
Effective PCU & 7.7\% & 68.9\% & 4.7\% & 7.1\% & 42.4\% & 51.5\% & 23.4\% \\
\hline
\textbf{Badania kodów równoległych} & \textbf{k1} & \textbf{k2} & \textbf{k3} & \textbf{k3a} & \textbf{k4} & \textbf{k4a} & \textbf{k5} \\
\hline
Przyspieszenie względem k3a & - & 0.05 & - & - & 0.34 & 0.43 & 2.17 \\
\hline
Efektywność [Acc/Cores] & - & 0.004 & - & - & 0.028 & 0.036 & 0.180 \\
\hline
\end{tabular}%
% }
\endgroup
\caption{Porównanie podstawowych parametrów wywołania dla kodów k1-k5.}
\end{table}
\begin{table}[h!]
\centering
\begingroup
% \resizebox{\textwidth}{!}{%
\begin{tabular}{|l|c|c|c|c|c|c|c|}
\hline
\textbf{Performance Cores} & \textbf{k1} & \textbf{k2} & \textbf{k3} & \textbf{k3a} & \textbf{k4} & \textbf{k4a} & \textbf{k5} \\
\hline
Retiring & 66.0\% & 83.5\% & 1.2\% & 0.0\% & 1.1\% & 4.5\% & 25.0\% \\
\hline
Front End Bound & 11.4\% & 13.8\% & 0.0\% & 3.1\% & 3.6\% & 20.0\% & 25.0\% \\
\hline
Back End Bound & 3.3\% & 1.7\% & 100.0\% & 0.0\% & 78.7\% & 58.0\% & 41.7\% \\
\hline
Memory Bound & 0.2\% & 0.2\% & 100.0\% & 0.0\% & 67.4\% & 42.6\% & 16.7\% \\
\hline
Core Bound & 3.1\% & 1.5\% & 0.0\% & 0.0\% & 11.4\% & 15.3\% & 25.0\% \\
\hline
\textbf{Efficient Cores} & \textbf{k1} & \textbf{k2} & \textbf{k3} & \textbf{k3a} & \textbf{k4} & \textbf{k4a} & \textbf{k5} \\
\hline
Retiring & 0.0\% & 42.3\% & 0.0\% & 0.0\% & 3.0\% & 0.7\% & 0.0\% \\
\hline
Front End Bound & 0.0\% & 5.0\% & 0.0\% & 0.0\% & 1.3\% & 3.8\% & 0.0\% \\
\hline
Back End Bound & 4.1\% & 11.5\% & 0.0\% & 0.0\% & 90.8\% & 34.8\% & 97.4\% \\
\hline
Memory Bound & 0.0\% & 0.2\% & 0.0\% & 0.0\% & 53.0\% & 34.8\% & 18.8\% \\
\hline
Core Bound & 4.1\% & 11.3\% & 0.0\% & 0.0\% & 37.8\% & 0.0\% & 78.6\% \\
\hline
\end{tabular} 
% }
\endgroup
\caption{Wyniki analizy przyczyn ograniczeń wydajności dla kodów k1-k5 z podziałem na Performance i Efficient Cores.}
\end{table}
\newpage
\subsection*{Porównanie działań kodów dla różnych zakresów liczb.}
W ramach tego eksperymentu zmierzyliśmy Ellapsed Time dla kodów k1-k5 dla trzech
różnych zakresów liczb: \\<2; max>, <max/2; max> oraz <2; max/2>, gdzie \texttt{max} to 10\textsuperscript{8}.
\begin{table}[h!]
\centering
\begingroup
% \resizebox{\textwidth}{!}{%
\begin{tabular}{|l|c|c|c|c|c|c|c|}
\hline
\textbf{Zakres liczb} & \textbf{k1} & \textbf{k2} & \textbf{k3} & \textbf{k3a} & \textbf{k4} & \textbf{k4a} & \textbf{k5} \\
\hline
<2; max> & 23.750 s & 2.394 s & 0.660 s & 0.130 s & 0.381 s & 0.296 s & 0.060 s \\
\hline
<max/2; max> & 16.671 s & 3.852 s & 0.303 s & 0.069 s & 0.279 s & 0.122 s & 0.038 s \\
\hline
<2; max/2> & 7.116 s & 2.450 s & 0.282 s & 0.065 s & 0.225 s & 0.117 s & 0.050 s \\
\hline

\end{tabular}%
% }
\endgroup
\caption{Porównanie prędkości działania dla różnych zakresów liczb w kodach k1-k5.}
\end{table}
\subsection*{Badanie wpływu różnych szeregowań dla czasu przetwarzania kodów k2, k4, k4a i k5.}

\begin{table}[h!]
\centering
\begingroup
% \resizebox{\textwidth}{!}{%
\begin{tabular}{|l|c|c|c|c|}
\hline
\textbf{Schedule} & \textbf{k2} & \textbf{k4} & \textbf{k4a} & \textbf{k5} \\
\hline
static & 6.339 s & 0.551 s & 0.541 s & 0.151 s \\
\hline
dynamic & 2.475 s & 0.296 s & 0.381 s & 0.080 s \\
\hline
guided &  2.394 s & 0.558 s & 0.727 s & 0.060 s \\
\hline

\end{tabular}%
% }
\endgroup
\caption{Wpływy szeregowań na czas przetwarzania kodów równoległych.}
\end{table}

\subsection*{Porównanie parametrów przetwarzania kodów k3 i k3a.}
\subsection*{Wyniki eksperymentu wyznaczania wielkości bloku dla kodów k3a i k5.}

Aby wyznaczyć optymalną wielkość bloku dla kodów k3a i k5 napisaliśmy eksperyment,
który wywoływał te kody dla standardowego zakres <2; max> z różnymi rozmiarami bloków
od 16B do ilości liczb w zakresie z krokiem x2. Poniżej znajduje się część tabeli wynikowej,
w której znajduje się nasze poszukiwane optimum (czyli w naszym przypadku 65536 B tj 64 kB).

\begin{table}[h!]
\centering
\begingroup
% \resizebox{\textwidth}{!}{%
\begin{tabular}{|l|c|c|}
\hline
\textbf{Blocksize} & \textbf{k3a} & \textbf{k5} \\
\hline
8192 & 0.229 s & 0.040 s \\
\hline
16384 & 0.145 s & 0.036 s \\
\hline
32768 & 0.112 s & 0.030 s \\
\hline
65536 & 0.107 s & 0.030 s \\
\hline
131072 & 0.121 s & 0.039 s \\
\hline
262144 & 0.126 s & 0.035 s \\
\hline
524288 & 0.128 s & 0.040 s \\
\hline


\end{tabular}%
% }
\endgroup
\caption{Wpływ rozmiaru bloku na czas przetwarzania kodów k3a i k5.}
\end{table}

\newpage
\section{Wnioski.}


\end{document}
